\documentclass[12pt,a4paper,oneside]{report}
% !TEX program = xelatex
\usepackage{fix-cm}
\usepackage{fontspec}
\setmainfont{Times New Roman}
\usepackage[vietnamese]{babel}
\usepackage[left=3.5cm,right=2cm,top=2.5cm,bottom=2.5cm]{geometry}
\usepackage{setspace}
\onehalfspacing
\usepackage{indentfirst}
\setlength{\parindent}{1cm}
\usepackage{graphicx,float,booktabs,longtable,tabularx,amsmath,enumitem}
\DeclareMathSizes{13}{12}{8}{6}
\usepackage[numbers,sort&compress]{natbib}
\usepackage{hyperref}
\hypersetup{colorlinks=true,linkcolor=black,citecolor=blue,urlcolor=blue}
\renewcommand{\contentsname}{MỤC LỤC}
\renewcommand{\listfigurename}{DANH MỤC HÌNH}
\renewcommand{\listtablename}{DANH MỤC BẢNG}
\renewcommand{\bibname}{TÀI LIỆU THAM KHẢO}
\graphicspath{{figures/}}
\begin{document}

%---------------- BIA ----------------
\hypersetup{pageanchor=false}
\begin{titlepage}
\fontsize{13pt}{18pt}\selectfont
\begin{center}
{\bfseries\Large ỦY BAN NHÂN DÂN TP. HỒ CHÍ MINH\\TRƯỜNG ĐẠI HỌC SÀI GÒN}
\vspace*{2cm}

{\bfseries\huge TIỂU LUẬN MÔN HỌC}

\vspace{0.5cm}
\textbf{PHƯƠNG PHÁP NGHIÊN CỨU KHOA HỌC CHO KHDL\\Mã MH: 858025}

\vspace*{2cm}
{\bfseries\huge PHÂN TÍCH MỐI LIÊN HỆ GIỮA HÀNH VI SỬ DỤNG ĐIỆN THOẠI TRƯỚC KHI NGỦ,
CHẤT LƯỢNG GIẤC NGỦ VÀ DỰ BÁO NỢ NGỦ BẰNG HỌC MÁY}

\vspace{2.5cm}
\begin{tabbing}
		\hspace*{2cm} {\large DANH SÁCH NHÓM:12} \\		 

		\hspace*{4cm} {\large 1. Họ và tên: Phan Trần Hữu Tấn} \ \ \ \ \ Lớp: DDU1231 \\ \ \ \ \ 
        \hspace*{4cm} {\large MSSV: 3123580042}\\
		
		\hspace*{4cm} {\large 2. Họ và tên: Trần Quốc Hoàng}\ \ \ \ \ Lớp: DDU1231 \\ \ \ \ \
        \hspace*{4cm} {\large MSSV: 3123580042}\\

        \hspace{2cm}\textbf{\large GIẢNG VIÊN GIẢNG DẠY:} ThS. Trương Hoàng Huy
	\end{tabbing}    

    \vspace*{0.5cm}
	
	\begin{center}
		{\bf Thành phố Hồ Chí Minh  \\  2026}
	\end{center}
\end{center}
\newpage
\hypersetup{pageanchor=true}

\pagenumbering{roman}
\chapter*{LỜI CAM ĐOAN}
Nhóm cam đoan báo cáo được thực hiện phục vụ học tập và nghiên cứu trong khuôn khổ môn học. Các nguồn tài liệu sử dụng được trích dẫn trong bài. Kết quả, bảng biểu và hình ảnh sẽ được xây dựng từ quá trình phân tích thực tế của nhóm; nhóm chịu trách nhiệm về tính trung thực của nội dung do mình thực hiện.
\clearpage

\chapter*{LỜI CẢM ƠN}
Nhóm xin chân thành cảm ơn giảng viên ThS. Trương Hoàng Huy đã hướng dẫn và cung cấp kiến thức trong môn Phương pháp nghiên cứu khoa học cho Khoa học dữ liệu. Do giới hạn về thời gian và kinh nghiệm, báo cáo có thể còn thiếu sót; nhóm mong nhận được góp ý để hoàn thiện.
\clearpage

\chapter*{TÓM TẮT}
Đề tài phân tích mối liên hệ giữa hành vi sử dụng điện thoại trước khi ngủ, các chỉ số giấc ngủ và mức độ nợ ngủ dựa trên bộ dữ liệu \texttt{bedtime\_screentime\_sleep\_debt.csv}. Quy trình dự kiến gồm kiểm tra chất lượng dữ liệu, phân tích khám phá dữ liệu, kiểm định sự khác biệt giữa nhóm, xây dựng các mô hình phân loại và giải thích kết quả. Các mô hình có thể gồm Logistic Regression, Random Forest và XGBoost; hiệu suất được đánh giá bằng Precision, Recall, F1-score, Macro-F1 và ma trận nhầm lẫn. Kết quả và kết luận cuối cùng sẽ được bổ sung sau khi nhóm hoàn tất thực nghiệm; báo cáo không giả định trước mô hình tốt nhất hoặc quan hệ nhân quả.
\\
\noindent\textbf{Từ khóa:} điện thoại trước khi ngủ, giấc ngủ, nợ ngủ, khoa học dữ liệu, học máy, SHAP.
\clearpage

\tableofcontents\clearpage
\listoffigures\clearpage
\listoftables\clearpage
\chapter*{DANH MỤC TỪ VIẾT TẮT}
\begin{longtable}{p{3cm}p{10cm}}
\toprule
EDA & Phân tích khám phá dữ liệu (Exploratory Data Analysis)\\
ML & Học máy (Machine Learning)\\
XAI & Trí tuệ nhân tạo có thể giải thích (Explainable AI)\\
SHAP & SHapley Additive exPlanations\\
RF & Random Forest\\
\bottomrule
\end{longtable}
\clearpage
\pagenumbering{arabic}

\chapter{GIỚI THIỆU ĐỀ TÀI}
\section{Lý do chọn đề tài}
Điện thoại thông minh được sử dụng rộng rãi cho học tập, công việc và giải trí. Việc sử dụng thiết bị vào buổi tối có thể liên hệ với thời điểm đi ngủ, độ trễ bắt đầu ngủ và trạng thái tỉnh táo vào ngày hôm sau. Mối liên hệ này có thể thay đổi theo thói quen sinh hoạt, loại ứng dụng, cài đặt màn hình, caffeine và đặc điểm cá nhân. Từ góc độ khoa học dữ liệu, bộ dữ liệu \texttt{bedtime\_screentime\_sleep\_debt.csv} cho phép nhóm khảo sát các mối liên hệ và đánh giá khả năng dự báo nợ ngủ bằng học máy.
\section{Mục tiêu nghiên cứu}
\subsection{Mục tiêu tổng quát}
Phân tích mối liên hệ giữa hành vi sử dụng điện thoại trước khi ngủ và các chỉ số liên quan đến giấc ngủ, đồng thời đánh giá mô hình học máy dự báo nhóm nợ ngủ.
\subsection{Mục tiêu cụ thể}
\begin{itemize}[leftmargin=1.2cm]
\item Kiểm tra chất lượng và mô tả bộ dữ liệu.
\item Phân tích tương quan giữa thời lượng dùng điện thoại và các chỉ số giấc ngủ.
\item So sánh nhóm nghề nghiệp, ứng dụng trước giờ ngủ và các kết quả liên quan.
\item Huấn luyện, so sánh mô hình phân loại \texttt{sleep\_debt\_category}.
\item Giải thích đóng góp của các đặc trưng và nêu hạn chế.
\end{itemize}
\section{Câu hỏi nghiên cứu}
\begin{enumerate}[leftmargin=1.2cm]
\item Thời lượng dùng điện thoại (\texttt{bedtime\_phone\_minutes}) liên hệ thế nào với độ trễ giấc ngủ và tổng thời gian ngủ?
\item Độ sáng màn hình và bộ lọc ánh sáng xanh liên hệ thế nào với tỷ lệ giấc ngủ sâu và REM?
\item Sự kết hợp giữa sử dụng điện thoại muộn và caffeine sau 17 giờ liên hệ thế nào với mệt mỏi ngày hôm sau?
\item Tỷ lệ nợ ngủ có khác nhau giữa các nhóm nghề nghiệp không?
\item Số lần báo thức lại có khác nhau giữa các nhóm ứng dụng trước giờ ngủ không?
\item Các mô hình học máy dự báo \texttt{sleep\_debt\_category} đạt hiệu suất như thế nào?
\item Những đặc trưng nào đóng góp nhiều nhất vào dự đoán của mô hình?
\end{enumerate}
\section{Đối tượng và phạm vi}
Đối tượng là các bản ghi trong tệp dữ liệu nêu trên. Nhóm cần ghi rõ nguồn, cách thu thập, thời điểm thu thập và cách lấy mẫu nếu thông tin được cung cấp. Không suy đoán những thông tin chưa được xác minh.
\section{Ý nghĩa và bố cục báo cáo}
Nghiên cứu minh họa việc kết hợp EDA, kiểm định thống kê, học máy và giải thích mô hình. Nếu dữ liệu là dữ liệu quan sát, kết quả không tự động chứng minh quan hệ nhân quả. Báo cáo gồm năm chương: giới thiệu; cơ sở lý thuyết; phân tích dữ liệu; mô hình học máy; kết luận và kiến nghị.

\chapter{CƠ SỞ LÝ THUYẾT VÀ NGHIÊN CỨU LIÊN QUAN}
\section{Chất lượng giấc ngủ và nợ ngủ}
Trong báo cáo, các khía cạnh liên quan đến giấc ngủ được khảo sát qua độ trễ đi vào giấc ngủ, tổng thời gian ngủ, tỷ lệ giấc ngủ sâu, tỷ lệ REM và điểm mệt mỏi ngày hôm sau. Cần trích dẫn nguồn khoa học khi giải thích ý nghĩa của các chỉ số. Phải xác minh định nghĩa và quy tắc gán nhãn của \texttt{sleep\_debt\_category}, bao gồm số lớp và ý nghĩa từng lớp.
\section{Hành vi sử dụng điện thoại và yếu tố liên quan}
Các biến thời lượng sử dụng điện thoại, ứng dụng chính, độ sáng và bộ lọc ánh sáng xanh mô tả các khía cạnh khác nhau của hành vi sử dụng thiết bị. Caffeine, hoạt động thể chất, tuổi, nghề nghiệp và chronotype có thể là các biến liên quan hoặc yếu tố gây nhiễu.
\section{Phương pháp thống kê}
Thống kê mô tả gồm trung bình, độ lệch chuẩn, trung vị, khoảng tứ phân vị, tần suất và tỷ lệ. Pearson đo liên hệ tuyến tính giữa hai biến số; Spearman có thể phù hợp hơn khi dữ liệu không đáp ứng giả định cần thiết. So sánh nhóm có thể dùng t-test/ANOVA hoặc kiểm định phi tham số tùy dữ liệu; biến phân loại có thể dùng Chi-square hoặc Fisher. Cần báo cáo giả thuyết, p-value, kích thước ảnh hưởng và xử lý so sánh nhiều lần khi phù hợp.
\section{Mô hình học máy}
\subsection{Logistic Regression}
Mô hình cơ sở có khả năng diễn giải tương đối thuận tiện. Với nhiều lớp, cần xác định chiến lược đa lớp và tiền xử lý phù hợp.
\subsection{Random Forest}
Random Forest kết hợp nhiều cây quyết định để mô hình hóa quan hệ phi tuyến. Feature importance không đồng nghĩa với tác động nhân quả \citep{breiman2001randomforest}.
\subsection{XGBoost}
XGBoost là phương pháp boosting dựa trên cây, thường được dùng với dữ liệu dạng bảng \citep{chen2016xgboost}. Siêu tham số cần được chọn bằng validation hoặc cross-validation.
\section{Đánh giá và giải thích mô hình}
Báo cáo Accuracy cùng Precision, Recall, F1-score theo lớp và Macro-F1. Ma trận nhầm lẫn cho biết các lớp bị dự đoán nhầm. SHAP hỗ trợ diễn giải đóng góp của đặc trưng vào dự đoán nhưng không chứng minh quan hệ nhân quả \citep{lundberg2017shap}.
\section{Khung và giả thuyết nghiên cứu}
Quy trình gồm kiểm tra dữ liệu, EDA và kiểm định, xây dựng mô hình, giải thích và thảo luận. Các giả thuyết có thể gồm: thời lượng dùng điện thoại liên hệ với độ trễ giấc ngủ; thời lượng dùng điện thoại liên hệ với tổng thời gian ngủ; caffeine và thời lượng điện thoại có thể có tương tác với mệt mỏi; phân bố nợ ngủ khác nhau theo nghề nghiệp; số lần báo thức lại khác nhau theo nhóm ứng dụng. Đây là giả thuyết cần kiểm tra, không phải kết luận có sẵn.

\chapter{DỮ LIỆU VÀ PHÂN TÍCH KHÁM PHÁ}
\section{Mô tả bộ dữ liệu}
Bộ dữ liệu có tên \texttt{bedtime\_screentime\_sleep\_debt.csv}. Nhóm cần nêu nguồn dữ liệu, đơn vị quan sát, cách thu thập, thời gian thu thập và giấy phép sử dụng nếu có.
\begin{table}[H]\centering
\caption{Thông tin tổng quan bộ dữ liệu}
\begin{tabularx}{0.9\textwidth}{lX}\toprule
Tên tệp & \texttt{bedtime\_screentime\_sleep\_debt.csv}\\
Số dòng & [Điền từ kết quả kiểm tra]\\
Số cột & [Điền từ kết quả kiểm tra]\\
Nguồn dữ liệu & [Điền nguồn đã xác minh]\\
\bottomrule\end{tabularx}\end{table}
\section{Kiểm tra chất lượng dữ liệu}
Kiểm tra giá trị thiếu, dòng trùng lặp, kiểu dữ liệu, giá trị ngoài phạm vi hợp lý và phân phối bất thường. Ghi số lượng trước/sau xử lý và lý do cho mỗi quyết định.
\begin{table}[H]\centering
\caption{Tổng hợp chất lượng dữ liệu}
\begin{tabularx}{0.9\textwidth}{Xrr}\toprule
Tiêu chí & Trước xử lý & Sau xử lý\\\midrule
Số dòng & [ ] & [ ]\\
Số ô thiếu & [ ] & [ ]\\
Số dòng trùng & [ ] & [ ]\\
Giá trị ngoài phạm vi & [ ] & [ ]\\\bottomrule
\end{tabularx}\end{table}
\section{Thống kê mô tả}
Trình bày số lượng, tỷ lệ, trung bình, độ lệch chuẩn, trung vị, khoảng tứ phân vị, min và max tùy kiểu biến. Không điền số liệu khi chưa chạy mã phân tích.
\section{Tương quan giữa thời lượng dùng điện thoại và giấc ngủ}
Khảo sát \texttt{bedtime\_phone\_minutes} với \texttt{sleep\_latency\_min} và \texttt{total\_sleep\_hours} bằng biểu đồ phân tán, hệ số Pearson hoặc Spearman tùy phân phối.
\begin{figure}[H]\centering
\fbox{\parbox[c][4.5cm][c]{0.82\textwidth}{\centering Chèn biểu đồ phân tán hoặc biểu đồ tương quan}}
\caption{Mối liên hệ giữa thời lượng dùng điện thoại và chỉ số giấc ngủ}\label{fig:phone-sleep}
\end{figure}
\noindent\textbf{Kết quả:} [Điền hệ số, khoảng tin cậy nếu có, p-value và diễn giải.]
\section{Độ sáng, bộ lọc ánh sáng xanh và pha ngủ}
So sánh \texttt{deep\_sleep\_pct} và \texttt{rem\_sleep\_pct} theo độ sáng màn hình và trạng thái \texttt{blue\_light\_filter\_active}. Cân nhắc yếu tố gây nhiễu và kích thước nhóm.
\noindent\textbf{Kết quả:} [Điền biểu đồ, kiểm định và diễn giải.]
\section{Sử dụng điện thoại, caffeine và mệt mỏi}
Có thể khảo sát biến tương tác bằng mô hình hồi quy phù hợp. Ví dụ minh họa:
\[
Y=\beta_0+\beta_1X_{\text{phone}}+\beta_2X_{\text{caffeine}}+
\beta_3(X_{\text{phone}}\times X_{\text{caffeine}})+\epsilon.
\]
Trong đó \(Y\) là điểm mệt mỏi ngày hôm sau. Chọn mô hình cuối cùng dựa trên kiểu dữ liệu và giả định.
\noindent\textbf{Kết quả:} [Điền hệ số, khoảng tin cậy, p-value và diễn giải.]
\section{So sánh theo nghề nghiệp và ứng dụng}
Lập bảng chéo tỷ lệ nợ ngủ theo nghề nghiệp. So sánh số lần hoãn báo thức theo ứng dụng dùng trước khi ngủ. Chọn kiểm định phù hợp và báo cáo kích thước ảnh hưởng khi có thể.
\noindent\textbf{Kết quả nghề nghiệp:} [Điền bảng và kết quả kiểm định.]
\\
\noindent\textbf{Kết quả ứng dụng:} [Điền thống kê và kết quả kiểm định.]
\section{Tóm tắt EDA}
Tóm tắt phát hiện cho từng câu hỏi nghiên cứu, bằng chứng thống kê và giới hạn. Không dùng kết quả từ tập kiểm tra để lựa chọn đặc trưng cho mô hình.

\chapter{XÂY DỰNG VÀ ĐÁNH GIÁ MÔ HÌNH HỌC MÁY}
\section{Xác định bài toán và kiểm tra nhãn}
Biến mục tiêu là \texttt{sleep\_debt\_category}. Xác minh số lớp, số mẫu từng lớp và quy tắc gán nhãn. Nếu nhãn được tạo trực tiếp từ \texttt{total\_sleep\_hours} hoặc một biến đầu vào khác, cần đánh giá nguy cơ rò rỉ nhãn trước khi đưa biến đó vào mô hình.
\section{Tiền xử lý và lựa chọn đặc trưng}
Mô tả cách xử lý giá trị thiếu, mã hóa biến phân loại, chuẩn hóa khi cần, và xử lý mất cân bằng lớp. Các bước học từ dữ liệu phải được fit chỉ trên tập huấn luyện trong từng fold. Nêu rõ biến nào được sử dụng và lý do loại bỏ các biến không phù hợp.
\section{Chia dữ liệu và đánh giá}
Có thể chia train/test có phân tầng theo nhãn nếu phù hợp. Dùng cross-validation trên tập huấn luyện để chọn siêu tham số; chỉ dùng tập kiểm tra cho đánh giá cuối cùng.
\begin{itemize}[leftmargin=1.2cm]
\item Tỷ lệ train/test: [Điền].
\item Số fold: [Điền].
\item Random seed: [Điền].
\item Xử lý mất cân bằng: [Điền].
\item Phiên bản thư viện: [Điền].
\end{itemize}
\section{Mô hình}
\subsection{Logistic Regression}
Mô hình cơ sở; mô tả mã hóa, chuẩn hóa, chiến lược đa lớp và tham số.
\subsection{Random Forest}
Ghi các tham số như số cây, độ sâu tối đa và cách chọn tham số.
\subsection{XGBoost}
Nếu sử dụng, ghi phiên bản, hàm mục tiêu và quá trình điều chỉnh tham số. Nếu không dùng được XGBoost, giải thích thuật toán thay thế.
\section{Chỉ số đánh giá}
Báo cáo Accuracy, Precision, Recall, F1-score theo lớp, Macro-F1 và ma trận nhầm lẫn. Với lớp hiếm, chú ý Recall của lớp nợ ngủ nghiêm trọng.
\begin{table}[H]\centering
\caption{So sánh hiệu suất mô hình}
\begin{tabularx}{\textwidth}{Xrrrr}\toprule
Mô hình & Accuracy & Macro-F1 & Recall lớp nghiêm trọng & Thời gian\\\midrule
Logistic Regression & [ ] & [ ] & [ ] & [ ]\\
Random Forest & [ ] & [ ] & [ ] & [ ]\\
XGBoost & [ ] & [ ] & [ ] & [ ]\\\bottomrule
\end{tabularx}\end{table}
\section{Ma trận nhầm lẫn}
\begin{figure}[H]\centering
\fbox{\parbox[c][4.5cm][c]{0.8\textwidth}{\centering Chèn ma trận nhầm lẫn của mô hình cuối cùng}}
\caption{Ma trận nhầm lẫn trên tập kiểm tra}\label{fig:confusion}
\end{figure}
Phân tích các lớp thường bị nhầm và các lỗi liên quan đến lớp nợ ngủ nghiêm trọng.
\section{Giải thích mô hình bằng Feature Importance và SHAP}
Có thể dùng Feature Importance và SHAP để khảo sát đóng góp của đặc trưng. Ghi rõ mô hình, tập dữ liệu và cách tổng hợp SHAP. Mức quan trọng không chứng minh quan hệ nhân quả.
\begin{figure}[H]\centering
\fbox{\parbox[c][4.5cm][c]{0.8\textwidth}{\centering Chèn biểu đồ Feature Importance hoặc SHAP}}
\caption{Các đặc trưng quan trọng theo mô hình được chọn}\label{fig:shap}
\end{figure}
\section{Thảo luận và trả lời câu hỏi nghiên cứu}
RQ6: [So sánh mô hình theo chỉ số thực nghiệm và tiêu chí chính đã chọn.]\\
RQ7: [Nêu đặc trưng nổi bật và giới hạn diễn giải.]\\
Không điền số liệu hoặc kết luận trước khi chạy mô hình.

\chapter{KẾT LUẬN VÀ KIẾN NGHỊ}
\section{Kết luận}
Tổng hợp kết quả thực nghiệm từ Chương 3 và Chương 4. Trả lời từng câu hỏi nghiên cứu, phân biệt kết quả mô tả, kiểm định và dự báo.
\begin{itemize}[leftmargin=1.2cm]
\item Mối liên hệ giữa thời lượng dùng điện thoại và chỉ số giấc ngủ: [Điền].
\item Kết quả về độ sáng và bộ lọc ánh sáng xanh: [Điền].
\item Kết quả về caffeine và mệt mỏi: [Điền].
\item Khác biệt theo nghề nghiệp và ứng dụng: [Điền].
\item Mô hình phù hợp nhất theo tiêu chí đã chọn: [Điền].
\item Đặc trưng nổi bật theo phân tích giải thích: [Điền].
\end{itemize}
\section{Đóng góp của nghiên cứu}
Nêu đóng góp thực tế trong kiểm tra dữ liệu, phân tích thống kê, xây dựng quy trình đánh giá mô hình và diễn giải dự đoán. Chỉ khẳng định những đóng góp có bằng chứng trong báo cáo.
\section{Khuyến nghị}
Khuyến nghị phải dựa trên kết quả và tài liệu tham khảo. Nếu dữ liệu chỉ cho thấy tương quan, dùng cách diễn đạt “có liên hệ với” thay vì “gây ra”. Không đề xuất ngưỡng sử dụng điện thoại hoặc caffeine cụ thể nếu chưa có bằng chứng phù hợp.
\section{Hạn chế}
Thảo luận nguồn gốc và tính đại diện của dữ liệu, cách đo lường biến, yếu tố gây nhiễu, mất cân bằng lớp, nguy cơ rò rỉ nhãn, lựa chọn mô hình và giới hạn suy luận nhân quả. Nếu dữ liệu là dữ liệu mô phỏng/tổng hợp, phải nêu rõ.
\section{Hướng nghiên cứu tiếp theo}
Có thể mở rộng bằng dữ liệu theo thời gian, kiểm định trên bộ dữ liệu độc lập, bổ sung biến gây nhiễu, thử nghiệm mô hình khác và đánh giá độ ổn định của phương pháp giải thích.


\bibliographystyle{plainnat}
{\raggedright\bibliography{references}}

\appendix
\chapter{Mô tả biến và tái lập thực nghiệm}
Bổ sung bảng mô tả đầy đủ các biến, nguồn dữ liệu, định nghĩa nhãn, phiên bản thư viện, random seed, cấu hình chia tập và siêu tham số. Chỉ ghi thông tin đã xác minh.
\end{document}
